% Lösungen Einheit 1
\einheitenkopf[le1][1 Wurzelziehen]{Einheit 1 von 4 · Wurzelziehen und Lösbarkeit}
\erg{12}{a) quadratisch \quad b) linear \quad c) quadratisch \quad d) quadratisch, denn x · x = x² \quad e) linear}
\erg{13}{a) positiv \quad b) negativ \quad c) null \quad d) positiv \quad e) negativ}
\erg{14}{a) $x_1$ = 4, $x_2$ = −4 \quad b) $x_1$ = 9, $x_2$ = −9 \quad c) $x_1$ = 2, $x_2$ = −2 \quad d) $x_1$ = 12, $x_2$ = −12 \quad e) $x_1$ = $\surd$10 $\approx$ 3,16, $x_2$ = −$\surd$10 $\approx$ −3,16 \quad f) Die Gleichung hat keine Lösung.}
\erg{15}{a) x² = 16; $x_1$ = 4, $x_2$ = −4 \quad b) x² = 64; $x_1$ = 8, $x_2$ = −8 \quad c) x² = 25; $x_1$ = 5, $x_2$ = −5 \quad d) x² = 4; $x_1$ = 2, $x_2$ = −2 \quad e) x² = −4; keine Lösung \quad f) x² = 0; x = 0}
\erg{16}{a) x − 2 = 3 oder x − 2 = −3; $x_1$ = 5, $x_2$ = −1 \quad b) $x_1$ = 6, $x_2$ = −4 \quad c) $x_1$ = 7, $x_2$ = 3 \quad d) $x_1$ = 4, $x_2$ = 2 \quad e) x + 2 = 5 oder x + 2 = −5; $x_1$ = 3, $x_2$ = −7 \quad f) x + 5 = 0; x = −5}
\erg{17}{a) 49 + 1 = 50, wA \quad b) 49 + 1 = 50, wA \quad c) (3 − 1)² = 4 $\neq$ 9, fA \quad d) (−2 + 5)² = 9, wA \quad e) 3 · 1 − 5 = −2 $\neq$ −8, fA}
\erg{18}{a) 2 Lösungen: $x_1$ = 1, $x_2$ = −3 \quad b) 2 Lösungen: $x_1$ = 0, $x_2$ = −2 \quad c) 1 Lösung: x = −1 \quad d) keine Lösung}
\erg{19}{a) 0 \quad b) jede positive Zahl, z. B. 9 \quad c) jede Zahl kleiner als 5, z. B. 1 \quad d) jede negative Zahl, z. B. −3 \quad e) z. B. (x − 1)² = −4}
\erg{20}{a) erste Aussage wahr (nur x = 6); zweite Aussage falsch: x + 6 = 4 oder x + 6 = −4, also $x_1$ = −2, $x_2$ = −10 \quad b) z. B. x² = −4}
\erg{21}{a) Die negative Lösung fehlt: $x_1$ = 8, $x_2$ = −8 \quad b) 9 wurde addiert statt subtrahiert: $x_1$ = −7, $x_2$ = −11 \quad c) Aus einer negativen Zahl gibt es keine Wurzel, (−9)² = 81: keine Lösung}
\erg{22}{a) 20 ist positiv; $\surd$20 und −$\surd$20 ergeben quadriert beide 20. \quad b) Ein Quadrat ist nie negativ. \quad c) Das Quadrat ist nur null, wenn x − 9 = 0 ist, also nur für x = 9.}
\erg{23}{a) x² = 121, Seite 11 m (−11 ist keine Länge) \quad b) x² = 1,44, Seite 1,2 m \quad c) r² = 500 : (10$\pi$) $\approx$ 15,92; r $\approx$ 3,99 cm}
